% !TeX root = ../../Book2.tex
% 纲要标题：## 附录 A　结、链环与低维拓扑
% 纲要位置：几何学计划纲要.md，第 1974 行。
% 正文状态：未撰写；本文件目前是框架。
% 补充范围与先修：三维分解采用紧、连通、可定向三维流形的明确版本，边界环面情形另核条件；使用第一卷基本群、覆盖和本附录结补空间、Dehn 填充。素分解、JSJ 与几何化的深层定理准确外引；几何模型接第三卷空间形式及几何群论附录，Ricci 流输入见第三卷第27章，不承担一般三维分类证明。

% 纲要内容开始
% 
% > 不承担一般三维分类；从两个实际可计算不变量进入更深结构。
% 
% > 补充写作范围与外部输入：补充紧、连通、可定向三维流形的素分解、JSJ 与几何化图景，边界环面情形明确条件；深层定理准确外引。依赖第一卷基本群、覆盖及本附录的结补空间和 Dehn 填充；几何模型接第三卷空间形式及几何群论附录，Ricci 流输入见第三卷第27章，不承担一般三维分类证明。
% 
% * **结图与补空间**
%   * 环境同位及 Reidemeister 操作
%   * Wirtinger 展示及结群
%   * 链环数的相交定义
% 
% * **Seifert 曲面与 Alexander 不变量**
%   * Seifert 曲面的构造
%   * 无限循环覆盖和 Alexander 模
%   * Seifert 矩阵的多项式计算
%   * 三叶结和八字结的核验
% 
% * **三维手术**
%   * Dehn 填充和斜率
%   * 手术图及边界变化
%   * 三维分类所需的额外定理
% 
% * **三维分解与几何化**
%   * 连通和与素分解
%   * 不可约性与不可压缩环面
%   * JSJ 分解的选定陈述
%   * 八种几何与几何化定理
%   * 结补空间、Dehn 填充与几何变化
% 
% 
% ---
% 
% 纲要内容结束

\chapter{结、链环与低维拓扑}
\label{g2:app:A}
\BookChapterNumber{g2:app:A}

待撰写

\section{结图与补空间}
\label{g2:sec:A01}
待撰写

\subsection{环境同位及 Reidemeister 操作}
\label{g2:subsec:A01:01}
待撰写

\subsection{Wirtinger 展示及结群}
\label{g2:subsec:A01:02}
待撰写

\subsection{链环数的相交定义}
\label{g2:subsec:A01:03}
待撰写

\section{Seifert 曲面与 Alexander 不变量}
\label{g2:sec:A02}
待撰写

\subsection{Seifert 曲面的构造}
\label{g2:subsec:A02:01}
待撰写

\subsection{无限循环覆盖和 Alexander 模}
\label{g2:subsec:A02:02}
待撰写

\subsection{Seifert 矩阵的多项式计算}
\label{g2:subsec:A02:03}
待撰写

\subsection{三叶结和八字结的核验}
\label{g2:subsec:A02:04}
待撰写

\section{三维手术}
\label{g2:sec:A03}
待撰写

\subsection{Dehn 填充和斜率}
\label{g2:subsec:A03:01}
待撰写

\subsection{手术图及边界变化}
\label{g2:subsec:A03:02}
待撰写

\subsection{三维分类所需的额外定理}
\label{g2:subsec:A03:03}
待撰写

\section{三维分解与几何化}
\label{g2:sec:A04}
待撰写

\subsection{连通和与素分解}
\label{g2:subsec:A04:01}
待撰写

\subsection{不可约性与不可压缩环面}
\label{g2:subsec:A04:02}
待撰写

\subsection{JSJ 分解的选定陈述}
\label{g2:subsec:A04:03}
待撰写

\subsection{八种几何与几何化定理}
\label{g2:subsec:A04:04}
待撰写

\subsection{结补空间、Dehn 填充与几何变化}
\label{g2:subsec:A04:05}
待撰写
